% !TeX program = xelatex
% !TeX encoding = UTF-8
\documentclass[12pt,a4paper,twocolumn]{article}

% ------------------------------------------------------------
% بسته‌های مورد نیاز
% ------------------------------------------------------------
\usepackage{amsmath,amssymb,amsfonts}
\usepackage{booktabs,array,multirow}
\usepackage{geometry}
\usepackage{graphicx}
\usepackage{xcolor}
\usepackage{enumitem}
\usepackage{caption}
\usepackage{subcaption}
\usepackage{hyperref}
\usepackage{algorithm}
\usepackage{algpseudocode}
\usepackage{float}
\usepackage{url}

% ------------------------------------------------------------
% تنظیمات صفحه
% ------------------------------------------------------------
\geometry{left=2cm,right=2cm,top=2.5cm,bottom=2.5cm}
\setlength{\columnsep}{0.8cm}

% ------------------------------------------------------------
% تنظیمات فونت فارسی (نیازمند XeLaTeX)
% ------------------------------------------------------------
\usepackage{xepersian}
\settextfont[Scale=1.05]{B Nazanin}
\setlatintextfont[Scale=0.9]{Times New Roman}

% ------------------------------------------------------------
% تنظیمات عنوان
% ------------------------------------------------------------
\title{%
	\large
	موتور پردازش اقلیم‌شناسی: \\[0.3cm]
	\normalsize
	چارچوبی مقیاس‌پذیر برای انتخاب توزیع احتمالاتی و تحلیل مقادیر حدی در داده‌های اقلیمی ایستگاهی
}
\author{%
	امین‌فضل کاظمی \\[0.2cm]
	\small
	\texttt{aminfazlkazemi@gmail.com} \\[0.1cm]
	\small
	گروه فیزیک فضا، دانشکده فیزیک، دانشگاه تهران، تهران، ایران
}
\date{}

% ------------------------------------------------------------
% شروع سند
% ------------------------------------------------------------
\begin{document}

\maketitle

% ============================================================
% چکیده
% ============================================================
\begin{abstract}
\noindent
تحلیل آماری داده‌های اقلیمی در مقیاس بزرگ با چالش‌های مضاعف آماری و محاسباتی مواجه است. از یک سو، شکل توزیع داده‌ها در نقاط مختلف شبکه یکسان نیست و فرض یک توزیع ثابت می‌تواند به برآوردهای نادرست از احتمال رویدادها منجر شود. از سوی دیگر، برازش همزمان چندین مدل آماری برای صدها هزار ایستگاه و هزاران روز، نیازمند زیرساختی کارآمد از نظر حافظه، زمان پردازش و عملیات ورودی/خروجی است. در این پژوهش، «موتور پردازش اقلیم‌شناسی» به‌عنوان یک چارچوب متن‌باز و ماژولار معرفی می‌شود که چهار خانواده توزیعی شامل نرمال، نرمال چوله، مخلوط گاوسی دومؤلفه‌ای و گامای انتقال‌یافته را در یک گردش‌کار واحد برای مدل‌سازی داده‌های خام پنجره‌ای مقایسه می‌کند. انتخاب مدل با معیار اطلاعاتی آکائیکه تصحیح‌شده (AICc) انجام می‌شود و برای نمونه‌های کوچک، حداقل تعداد مشاهدات موردنیاز برای هر مدل به‌صورت جداگانه تعریف شده است. به‌طور جداگانه، توزیع مقادیر حدی تعمیم‌یافته (GEV) بر روی بیشینه‌های سالانه‌ی پنجره برازش می‌شود و در فرآیند انتخاب مدل داده‌های خام با AICc وارد نمی‌شود. برای مدیریت داده‌های حجیم، سامانه از پردازش بلوکی، کش‌سازی چنداندازه‌ی دیسک و سیستم checkpoint برای ازسرگیری خودکار استفاده می‌کند. پیاده‌سازی با پایتون و کتابخانه‌های Numba، Xarray و Zarr انجام شده است. ارزیابی روی مجموعه‌داده‌ای شامل ۳۳۸٬۶۲۷ ایستگاه و دوره‌ی ۳۱ ساله (۱۳۶۹–۱۳۹۹) نشان می‌دهد که سامانه با سرعت میانگین ۷٫۵ ایستگاه در ثانیه و اوج حافظه ۳٫۵۶ گیگابایت، کل پردازش را در حدود ۱۲٫۵ ساعت تکمیل می‌کند. سیستم کش، زمان اجرای مرتبط با ورودی/خروجی را حدود ۶۰\% کاهش می‌دهد. نتایج انتخاب مدل برای متغیر Tmax نشان می‌دهد که ۴۲.۳\% از برازش‌های معتبر به توزیع نرمال، ۲۸.۷\% به نرمال چوله، ۱۵.۲\% به مخلوط گاوسی دومؤلفه‌ای و ۵.۴\% به گامای انتقال‌یافته اختصاص دارد (۸.۴\% موارد فاقد برازش معتبر بودند). تحلیل GEV نیز گرایش به رفتار دم سنگین را در بخش قابل‌توجهی از ایستگاه‌ها نشان می‌دهد، هرچند عدم‌قطعیت برآوردها قابل‌توجه است. اعتبارسنجی با داده‌های مصنوعی نشان داد که چارچوب پیشنهادی توانایی مناسبی در شناسایی توزیع مولد دارد و بازتوانی پارامترهای GMM نیز خطاهای نسبی محدودی نشان داد. کد منبع و مستندات در مخزن عمومی GitHub در دسترس است.
\end{abstract}

\noindent
\textbf{واژه‌های کلیدی:} اقلیم‌شناسی، برازش توزیع، انتخاب مدل، AICc، GEV، مدل مخلوط، EM، Zarr، پردازش بلوکی، بازتولیدپذیری.

% ============================================================
% ۱. مقدمه
% ============================================================
\section{مقدمه}

تغییرات اقلیمی و افزایش فراوانی رویدادهای فرین، نیاز به تحلیل دقیق داده‌های اقلیمی را بیش از پیش آشکار ساخته است. یکی از روش‌های بنیادین در این حوزه، برازش توزیع‌های احتمالاتی به داده‌های سری‌زمانی است که امکان برآورد دوره‌های بازگشت، شناسایی روندهای تغییرپذیری و شبیه‌سازی داده‌های مصنوعی را فراهم می‌کند \cite{wilks2011}. با این حال، داده‌های اقلیمی الزاماً از یک شکل توزیعی واحد پیروی نمی‌کنند. برخی سری‌ها تقریباً متقارن هستند، برخی چولگی دارند، برخی ممکن است نشانه‌هایی از چند رژیم آماری داشته باشند، و برخی نیز رفتار حدی از خود نشان می‌دهند \cite{katz2002}. بنابراین، تحمیل یک توزیع ثابت به تمام نقاط شبکه می‌تواند منجر به برآوردهای نادرست از احتمال‌ها، صدک‌ها و رویدادهای حدی شود.

از سوی دیگر، حجم عظیم داده‌ها (صدها هزار ایستگاه در بازه‌های زمانی چندین‌دهه)، تنوع فرمت‌های داده، نیاز به پردازش تکراری و محدودیت‌های حافظه و زمان پردازش، چالش‌های عمده‌ای در پیاده‌سازی عملیاتی روش‌های آماری ایجاد می‌کنند. مقایسه همزمان چندین مدل آماری برای هر ایستگاه و هر روز از سال، تعداد عملیات را به سطحی می‌رساند که بدون مدیریت هوشمند حافظه و ورودی/خروجی، عملاً غیرممکن است.

هدف این مطالعه، توسعه و ارزیابی یک چارچوب محاسباتی مقیاس‌پذیر برای انتخاب توزیع احتمالاتی مناسب در داده‌های اقلیمی است. نوآوری‌های اصلی این پژوهش عبارتند از: (۱) معماری افزونه‌ای برای مدل‌های آماری که امکان افزودن توزیع‌های جدید را بدون تغییر در هسته فراهم می‌کند، (۲) مقایسه‌ی هم‌زمان چهار خانواده‌ی توزیعی شامل نرمال، نرمال چوله، مخلوط گاوسی دومؤلفه‌ای و گامای انتقال‌یافته برای داده‌های خام، و تحلیل جداگانه‌ی GEV برای داده‌های حدی، (۳) استفاده از معیار AICc با قانون جریمه‌ی سادگی برای انتخاب مدل بهینه در هر روز و هر ایستگاه، (۴) پیاده‌سازی مدل مخلوط گاوسی دومؤلفه‌ای با الگوریتم EM و شاخص‌های تشخیص جدایش (Ashman's D، ضریب همپوشانی و شاخص جدایش وزنی پیشنهادی)، و (۵) پردازش بلوکی، کش‌سازی هوشمند و سیستم checkpoint برای مدیریت داده‌های حجیم و افزایش بازتولیدپذیری.

ساختار مقاله به‌این‌صورت است: در بخش ۲، داده‌ها و طراحی مطالعه تشریح می‌شود. بخش ۳ به معرفی چارچوب آماری شامل چهار توزیع برای داده‌های خام و تحلیل جداگانه‌ی GEV اختصاص دارد. بخش ۴ معماری نرم‌افزار و پیاده‌سازی را توضیح می‌دهد. نتایج تجربی در بخش ۵ ارائه و در بخش ۶ بحث می‌شوند. بخش ۷ به بازتولیدپذیری و کنترل کیفیت می‌پردازد و بخش ۸ نتیجه‌گیری و کارهای آینده را بیان می‌کند.

% ============================================================
% ۲. داده‌ها و طراحی مطالعه
% ============================================================
\section{داده‌ها و طراحی مطالعه}

\subsection{مجموعه‌داده}
برای ارزیابی سامانه، از داده‌های روزانهٔ دمای کمینه، میانگین و بیشینه (Tmin، Tmean، Tmax) مربوط به ۳۳۸٬۶۲۷ ایستگاه در دورهٔ ۳۱ ساله (۱۳۶۹ تا ۱۳۹۹ خورشیدی) استفاده شده است. داده‌ها از سازمان هواشناسی کشور و مرکز ملی خشکسالی تهیه شده و به‌صورت فایل‌های Zarr ماهانه با ساختار (time, point) ذخیره شده‌اند. جدول \ref{tab:dataset} مشخصات مجموعه‌داده را خلاصه می‌کند.

\begin{table}[h]
\centering
\caption{مشخصات مجموعه‌داده‌ی مورد استفاده}
\label{tab:dataset}
\begin{tabular}{l r}
\toprule
\textbf{ویژگی} & \textbf{مقدار} \\
\midrule
تعداد ایستگاه‌ها & ۳۳۸٬۶۲۷ \\
دوره‌ی زمانی & ۱۳۶۹–۱۳۹۹ (۳۱ سال) \\
تفکیک زمانی & روزانه \\
متغیرها & Tmin, Tmean, Tmax \\
تعداد روزهای سال & ۳۶۶ (شامل سال کبیسه) \\
قالب داده & Zarr v2 با فشرده‌سازی Zstd \\
واحد پردازش & پنجره‌ی پنج‌روزه \\
منبع داده & سازمان هواشناسی کشور، مرکز ملی خشکسالی \\
\bottomrule
\end{tabular}
\end{table}

مدیریت داده‌های نامعتبر به‌صورت زیر انجام می‌شود: روزهایی که داده ندارند با مقدار NaN مشخص می‌شوند. برای هر ایستگاه و هر روز، اگر تعداد مشاهدات معتبر در پنجره‌ی پنج‌روزه کمتر از آستانه‌ی تعیین‌شده برای هر مدل باشد، آن روز از فرایند برازش حذف می‌شود. سال‌های کبیسه با نگهداری روز ۲۹ اسفند (۳۶۶-امین روز سال) و استفاده از اندیس مدولار ۳۶۶ برای پنجره‌ها مدیریت شده‌اند. حدود ۸۲\% از ایستگاه‌ها داده‌ی کامل ۳۱ ساله دارند و میانگین نرخ داده‌ی گم‌شده حدود ۴٫۷\% است.

\subsection{واحد محاسباتی و پنجره‌ی زمانی}
برای هر ایستگاه و هر روز از سال، پنجره‌ای به‌طول ۵ روز شامل روز هدف و دو روز قبل و بعد از آن (در مجموع ۵ روز) از تمام سال‌های دوره استخراج می‌شود. برای روزهای ابتدا و انتهای سال، پنجره با استفاده از اندیس مدولار ۳۶۶ پیاده‌سازی می‌شود تا پیوستگی سالانه حفظ شود. این روش برای هر روز سال، حداکثر ۱۵۵ مشاهده (۳۱ سال × ۵ روز) فراهم می‌کند. پنجره‌ی پنج‌روزه به‌منظور ایجاد سازش میان همگنی فصلی و افزایش حجم نمونه انتخاب شده است \cite{burnham2002}.

% ============================================================
% ۳. چارچوب آماری
% ============================================================
\section{چارچوب آماری}

در این بخش، چهار توزیع آماری برای مدل‌سازی داده‌های خام پنجره‌ای معرفی می‌شوند. جدول \ref{tab:distributions} نقش هر یک را در پوشش رفتارهای آماری مختلف نشان می‌دهد.

\begin{table}[h]
\centering
\caption{خانواده‌های توزیعی انتخاب‌شده برای داده‌های خام و نقش هر یک}
\label{tab:distributions}
\begin{tabular}{l c p{6cm}}
\toprule
\textbf{توزیع} & \textbf{$k$} & \textbf{کاربرد مفهومی} \\
\midrule
نرمال & ۲ & مدل پایه برای داده‌های تقریباً متقارن \\
نرمال چوله & ۳ & داده‌های نامتقارن با یک مد \\
مخلوط گاوسی دومؤلفه‌ای & ۵ & تشخیص/مدل‌سازی دو مؤلفه یا دو رژیم \\
گامای انتقال‌یافته & ۳ & مدل انعطاف‌پذیر برای داده‌های چوله با کران پایین \\
\bottomrule
\end{tabular}
\end{table}

\subsection{توزیع نرمال}
تابع چگالی احتمال توزیع نرمال به‌صورت زیر تعریف می‌شود:
\begin{equation}
f(x) = \frac{1}{\sigma\sqrt{2\pi}}\exp\left( - \frac{(x - \mu)^{2}}{2\sigma^{2}} \right)
\end{equation}
که در آن $\mu$ میانگین و $\sigma>0$ انحراف معیار است. برآورد پارامترها با روش بیشینه‌ی درست‌نمایی (MLE) انجام می‌شود و به‌عنوان مدل مرجع در سامانه عمل می‌کند. حداقل تعداد مشاهدات موردنیاز برای این مدل $n_{\min}=5$ در نظر گرفته شده است.

\subsection{توزیع نرمال چوله}
نرمال چوله با افزودن پارامتر شکل $\alpha$ به مدل نرمال، امکان نمایش عدم تقارن را فراهم می‌کند \cite{azzalini1985}:
\begin{equation}
f(x) = \frac{2}{\sigma}\phi\left( \frac{x - \mu}{\sigma} \right)\Phi\left( \alpha\frac{x - \mu}{\sigma} \right)
\end{equation}
که در آن $\phi$ و $\Phi$ به‌ترتیب تابع چگالی و تابع توزیع تجمعی نرمال استاندارد هستند. پارامترها عبارت‌اند از: $\mu$ (مکان)، $\sigma>0$ (مقیاس) و $\alpha$ (شکل، کنترل‌کننده‌ی چولگی). برآورد اولیه با روش گشتاورها انجام می‌شود \cite{hosking1990}. مراحل الگوریتم به‌صورت زیر است:
\begin{enumerate}
\item محاسبه‌ی گشتاورهای نمونه ($\bar{x}, s^2, \gamma_1$)
\item برآورد اولیه‌ی $\alpha$ از روی چولگی با استفاده از رابطه‌ی $\gamma_1 = \frac{4-\pi}{2} \left( \frac{\delta}{\sqrt{1-\delta^2}} \right)^3$ که $\delta = \alpha/\sqrt{1+\alpha^2}$
\item برآورد $\mu$ و $\sigma$ از میانگین و واریانس
\item بهینه‌سازی عددی log-likelihood با روش L-BFGS-B برای بهبود برازش
\item بررسی همگرایی؛ در صورت عدم همگرایی، پرچم \texttt{NO\_CONVERGENCE} ثبت می‌شود.
\end{enumerate}
حداقل تعداد مشاهدات برای این مدل $n_{\min}=7$ در نظر گرفته شده است.

\subsection{مدل مخلوط گاوسی دومؤلفه‌ای و الگوریتم EM}
مدل مخلوط گاوسی دومؤلفه‌ای ترکیبی از دو توزیع نرمال با وزن‌های متفاوت است \cite{mclachlan2000}:
\begin{equation}
f(x) = w_{1}\mathcal{N}(x;\mu_{1},\sigma_{1}^{2}) + w_{2}\mathcal{N}(x;\mu_{2},\sigma_{2}^{2})
\end{equation}
که در آن $w_1+w_2=1$، $0\le w_i\le 1$، و $\mu_1,\mu_2$ و $\sigma_1,\sigma_2$ به‌ترتیب میانگین‌ها و انحراف معیارهای دو مؤلفه هستند.

در این مطالعه، منظور از «مدل دوکوهانه» در سطح برازش، یک مدل مخلوط گاوسی دومؤلفه‌ای است. با این حال، وجود دو مؤلفه در مدل لزوماً به معنای دوکوهانه بودن تابع چگالی نهایی نیست؛ زیرا بسته به وزن‌ها، فاصله‌ی میانگین‌ها و واریانس مؤلفه‌ها، چگالی مخلوط می‌تواند همچنان تک‌مدی باشد. بنابراین، دوکوهانه بودن به‌عنوان یک ویژگی توصیفی مستقل از صرف انتخاب مدل بررسی می‌شود.

برای تشخیص دوکوهانه‌ای، سه شاخص محاسبه می‌شوند:
\begin{itemize}
\item \textbf{شاخص جدایش اشمن (Ashman's D):} $D = \dfrac{|\mu_{2} - \mu_{1}|}{\sqrt{(\sigma_{1}^{2} + \sigma_{2}^{2})/2}}$ که $D>2$ نشان‌دهندهٔ جدایش واضح است.
\item \textbf{ضریب همپوشانی (OVL):} $OVL = \int \min\{f_{1}(x),f_{2}(x)\} dx$ که با انتگرال‌گیری عددی محاسبه می‌شود.
\item \textbf{شاخص جدایش وزنی (پیشنهادی در این مطالعه):}
\[
\beta = \dfrac{w_{1}w_{2}(\mu_{2} - \mu_{1})^{2}}{\sigma_{1}\sigma_{2}}.
\]
این کمیت در این مطالعه به‌عنوان یک شاخص اکتشافی و بدون‌بُعد برای ترکیب وزن مؤلفه‌ها، فاصله‌ی میانگین‌ها و پراکندگی درونی تعریف شده است. برخلاف Ashman's D، برای $\beta$ در این مقاله آستانه‌ی جهانی یا معیار تصمیم‌گیری قطعی پیشنهاد نمی‌شود و مقدار آن صرفاً برای مقایسه‌ی نسبی شدت جدایش گزارش می‌شود.
\end{itemize}

برآورد پارامترها با الگوریتم EM انجام می‌شود. برای کاهش حساسیت به مقداردهی اولیه، سه مقداردهی اولیه‌ی مختلف به‌کار گرفته می‌شود:
\begin{enumerate}
\item $\mu_1^{(0)} = \bar{x} - s$، $\mu_2^{(0)} = \bar{x} + s$، $\sigma_1^{(0)} = \sigma_2^{(0)} = s/2$، $w_1^{(0)}=0.5$
\item $\mu_1^{(0)} = \bar{x} - 0.5s$، $\mu_2^{(0)} = \bar{x} + 0.5s$، $\sigma_1^{(0)} = \sigma_2^{(0)} = s/2$، $w_1^{(0)}=0.5$
\item $\mu_1^{(0)} = \text{quantile}(0.25)$، $\mu_2^{(0)} = \text{quantile}(0.75)$، $\sigma_1^{(0)} = s/2$، $\sigma_2^{(0)} = s/2$، $w_1^{(0)}=0.5$
\end{enumerate}
در مرحلهٔ E، مسئولیت هر مشاهده محاسبه می‌شود. برای جلوگیری از underflow در محاسبه‌ی مسئولیت‌ها، به‌جای جمع مستقیم چگالی‌ها، از نمایش لگاریتمی و عملگر log-sum-exp استفاده شد. این روش از افزودن ثابت عددی به مخرج جلوگیری کرده و پایداری محاسبات را در نواحی دم توزیع افزایش می‌دهد. در مرحلهٔ M، پارامترها به‌روزرسانی می‌شوند:
\begin{align}
w_1 &= \frac{1}{n}\sum_{i=1}^{n} \gamma_{i1}, \quad w_2 = 1-w_1 \\
\mu_1 &= \frac{\sum_i \gamma_{i1} x_i}{\sum_i \gamma_{i1}}, \quad
\mu_2 = \frac{\sum_i (1-\gamma_{i1}) x_i}{\sum_i (1-\gamma_{i1})} \\
\sigma_1^2 &= \frac{\sum_i \gamma_{i1} (x_i-\mu_1)^2}{\sum_i \gamma_{i1}}, \quad
\sigma_2^2 = \frac{\sum_i (1-\gamma_{i1}) (x_i-\mu_2)^2}{\sum_i (1-\gamma_{i1})}
\end{align}
الگوریتم تا زمانی ادامه می‌یابد که تغییر نسبی log-likelihood در دو تکرار متوالی کمتر از $10^{-6}$ باشد یا به حداکثر ۱۰۰۰ تکرار برسد. از میان سه initialization، مدلی با بیشترین log-likelihood انتخاب می‌شود. برای جلوگیری از degeneracy، محدودیت‌های $\sigma_i > 10^{-3}$ و $w_i > 10^{-3}$ اعمال می‌شوند. همچنین پس از پایان الگوریتم، برای رفع مسئله‌ی label switching، مؤلفه‌ها بر اساس $\mu_1 < \mu_2$ مرتب می‌شوند. حداقل تعداد مشاهدات برای این مدل $n_{\min}=15$ در نظر گرفته شده است.

\subsection{توزیع گامای انتقال‌یافته}
گامای انتقال‌یافته (Shifted Gamma) یک توزیع گامای سه‌پارامتری است \cite{hosking1990}:
\begin{equation}
f(x) = \frac{1}{\Gamma(\alpha)\beta^{\alpha}}(x-\gamma)^{\alpha-1}\exp\left(-\frac{x-\gamma}{\beta}\right)
\end{equation}
که در آن $\alpha>0$ پارامتر شکل، $\beta>0$ پارامتر مقیاس و $\gamma$ پارامتر مکان (حد پایین) است و $x>\gamma$. برآورد پارامترها با روش گشتاورها بر اساس روابط زیر انجام می‌شود:
\begin{align}
\alpha &= \frac{4}{\gamma_1^2}, \quad \beta = \frac{s}{\sqrt{\alpha}}, \quad \gamma = \bar{x} - \alpha\beta
\end{align}
که در آن $\gamma_1$ چولگی نمونه و $s$ انحراف معیار نمونه است. در مواردی که $|\gamma_1|$ بسیار کوچک باشد (نزدیک به صفر)، برازش این توزیع به‌عنوان حالت نزدیک به نرمال مدیریت شده و از تولید پارامترهای ناپایدار جلوگیری می‌شود. حداقل تعداد مشاهدات برای این مدل $n_{\min}=8$ در نظر گرفته شده است. توجه شود که این پیاده‌سازی نسخه‌ی عمومی تمام حالت‌های Pearson Type III را پوشش نمی‌دهد و صرفاً برای داده‌های با چولگی مثبت پارامتری‌سازی شده است.

\subsection{انتخاب مدل با AICc و BIC}
پس از برازش چهار توزیع بر روی داده‌های خام پنجره، سامانه با استفاده از معیارهای اطلاعاتی، مدل بهینه را انتخاب می‌کند \cite{burnham2002}:
\begin{align}
\text{AIC} &= 2k - 2\ln(\hat{L}) \\
\text{AICc} &= \text{AIC} + \frac{2k(k+1)}{n-k-1} \\
\text{BIC} &= k\ln(n) - 2\ln(\hat{L})
\end{align}
که در آن $k$ تعداد پارامترها، $n$ حجم نمونه و $\hat{L}$ بیشینه‌ی درست‌نمایی است. توجه می‌شود که AICc تنها زمانی محاسبه می‌شود که $n > k+1$؛ در غیر این صورت، مدل برای رقابت AICc معتبر شناخته نشده و از فرآیند انتخاب حذف می‌شود. مدلی که کمترین AICc را دارد انتخاب می‌شود. اگر اختلاف AICc بین دو مدل کمتر از ۲ باشد، مدل با تعداد پارامتر کمتر (مدل ساده‌تر) ترجیح داده می‌شود. قاعده‌ی انتخاب به‌صورت $\Delta AICc = AICc_i - \min(AICc)$ و در صورت $\Delta AICc < 2$، انتخاب مدل با کمترین $k$ تعریف می‌شود.

\subsection{تحلیل جداگانه‌ی GEV برای مقادیر حدی}
به‌دلیل تفاوت بنیادی در ماهیت داده‌ها، توزیع GEV به‌صورت جداگانه برای تحلیل رفتار حدی در نظر گرفته شده است و در فرآیند انتخاب مدل با AICc برای داده‌های خام شرکت داده نمی‌شود. توزیع GEV برای مدل‌سازی بیشینه‌های بلوکی طراحی شده است \cite{coles2001}:
\begin{equation}
f(x) = \frac{1}{\sigma}\left( 1 + \xi\frac{x - \mu}{\sigma} \right)^{- \left( \frac{1}{\xi} + 1 \right)}\exp\left( - \left( 1 + \xi\frac{x - \mu}{\sigma} \right)^{- \frac{1}{\xi}} \right)
\end{equation}
که در آن $\mu$ پارامتر مکان، $\sigma>0$ پارامتر مقیاس و $\xi$ پارامتر شکل است. $\xi>0$ (فرشه، دم سنگین)، $\xi<0$ (ویبول، دم محدود) و $\xi=0$ (گامبل، دم نمایی).

برای هر ایستگاه و هر روز تقویمی ($d$)، بیشینه‌ی سالانه‌ی پنجره‌ی پنج‌روزه حول روز $d$ به‌صورت زیر تعریف می‌شود:
\begin{equation}
M_y(d) = \max_{j \in \{-2,-1,0,1,2\}} X_{y,\, d+j}
\end{equation}
سپس دنباله‌ی $\{M_{1369}(d), M_{1370}(d), \dots, M_{1399}(d)\}$ تشکیل می‌شود که در حالت داده‌ی کامل شامل ۳۱ مقدار است. این دنباله نمایانگر بیشینه‌ی سالانه‌ی پنجره‌ی پنج‌روزه حول روز تقویمی $d$ است و با «بیشینه‌ی سالانه‌ی کل سال» در ادبیات کلاسیک مقادیر حدی تفاوت دارد. بنابراین، تحلیل GEV حاضر یک تحلیل مقادیر حدی فصلی/تقویمی مبتنی بر بلوک‌های پنج‌روزه است، نه تحلیل استاندارد annual maximum series بر مبنای کل روزهای سال. در پیاده‌سازی، از تابع \texttt{genextreme.fit} کتابخانهٔ SciPy با قرارداد پارامتر $c=-\xi$ استفاده شده است. حداقل تعداد مشاهدات برای این مدل $n_{\min}=10$ در نظر گرفته شده است.

پس از برازش، دوره‌های بازگشت با استفاده از فرمول زیر محاسبه می‌شوند:
\begin{equation}
z_T = \mu + \frac{\sigma}{\xi}\left[ \left(-\log(1-1/T)\right)^{-\xi} - 1 \right]
\end{equation}
که در آن $T$ دوره‌ی بازگشت (بر حسب سال) است. برای $\xi \to 0$، حالت Gumbel به‌صورت $z_T = \mu - \sigma \log(-\log(1-1/T))$ به‌کار می‌رود.

\textbf{توجه:} برازش‌های GEV برای روزهای مختلف سال به دلیل هم‌پوشانی پنجره‌ها مستقل فرض نمی‌شوند و خلاصه‌های شبکه‌ای از پارامترهای GEV نباید به‌عنوان ۳۶۶ نمونه مستقل تفسیر شوند. همچنین، از آنجا که سری GEV در هر روز تقویمی حداکثر شامل ۳۱ مقدار سالانه است، برآورد دوره‌های بازگشت بزرگ‌تر از طول مؤثر سری به‌طور قابل‌توجهی به پارامتر شکل و عدم‌قطعیت برازش حساس است. بنابراین، مقادیر $T=50$ سال در این مطالعه باید به‌عنوان برآوردهای برون‌یابی‌شده و نه مشاهدات تجربی مستقیم تفسیر شوند. به‌عنوان مثال، برای یک ایستگاه نمونه با $\mu=25^\circ C$، $\sigma=2^\circ C$ و $\xi=0.1$، دوره‌های بازگشت $z_{10}=30.2^\circ C$، $z_{20}=32.1^\circ C$ و $z_{50}=34.5^\circ C$ برآورد شدند (اعداد با دقت یک رقم اعشار گزارش شده‌اند).

% ============================================================
% ۴. معماری نرم‌افزار و پیاده‌سازی
% ============================================================
\section{معماری نرم‌افزار و پیاده‌سازی}

سامانه با زبان پایتون و بر پایهٔ چهار اصل معماری افزونه‌محور، گردش‌کار تکرارپذیر، جداسازی وظایف و پیکربندی متمرکز طراحی شده است. کد منبع در مخزن عمومی GitHub \cite{github} در دسترس است.

\subsection{ساختار مخزن}
مخزن به‌صورت زیر سازمان‌دهی شده است:
\begin{itemize}
\item \texttt{core/}: هستهٔ سامانه شامل موتور افزونه‌ها، رابط داده، schema ذخیره‌سازی، کنترل کیفیت و عدم‌قطعیت.
\item \texttt{plugins/distributions/}: پیاده‌سازی چهار توزیع آماری به‌صورت افزونه.
\item \texttt{notebooks/}: آموزش‌های تعاملی Jupyter.
\item \texttt{tests/}: آزمون‌های واحد.
\item \texttt{config.yaml}: فایل پیکربندی اصلی.
\item \texttt{main.py}: نقطهٔ ورود اجرا.
\end{itemize}

\subsection{پردازش بلوکی و مدیریت حافظه}
برای مدیریت مجموعه‌داده‌های بزرگ‌تر از حافظه، داده‌ها به بلوک‌هایی با اندازهٔ قابل‌پیکربندی (پیش‌فرض ۲۰۰۰ ایستگاه) تقسیم می‌شوند. هر بلوک به‌طور مستقل بارگذاری، پردازش و ذخیره می‌شود. آرایه‌های داده با دقت float32 ذخیره می‌شوند تا مصرف حافظه کاهش یابد، اما کلیه‌ی محاسبات آماری حساس (log-likelihood، EM، AICc) در دقت float64 انجام می‌شوند. پس از پردازش هر بلوک، garbage collection به‌صورت صریح فراخوانی می‌شود.

\subsection{سیستم کش‌سازی هوشمند}
سیستم کش دیسک با شناسایی خودکار بلوک‌های ۱۰۰۰، ۲۰۰۰ و ۵۰۰۰ ایستگاه، از بارگذاری مجدد داده‌ها جلوگیری می‌کند. کلید کش شامل block\_start، block\_size، year، month، var و sample\_hash است. فایل‌های کش با فشرده‌سازی gzip ذخیره می‌شوند و حجم کش تا ۲۰ گیگابایت مدیریت می‌شود. این سیستم زمان اجرای مرتبط با ورودی/خروجی را حدود ۶۰\% کاهش می‌دهد.

\subsection{مدیریت Checkpoint و ازسرگیری خودکار}
سیستم checkpoint پس از پردازش هر بلوک، آخرین نقطهٔ پردازش‌شده (block, station) را در فایل checkpoint.csv ذخیره می‌کند. در زمان راه‌اندازی مجدد، سامانه فایل Zarr خروجی را بررسی کرده و آخرین نقطهٔ معتبر را شناسایی می‌کند. این قابلیت امکان ازسرگیری خودکار از نقطهٔ شکست را بدون نیاز به پردازش مجدد بلوک‌های قبلی فراهم می‌کند. برای تضمین consistency، ابتدا داده‌ها در Zarr نهایی شده و سپس checkpoint به‌روزرسانی می‌شود.

\subsection{ساختار داده‌های خروجی}
خروجی به‌صورت Zarr با ابعاد (day: 366, point: N\_stations) و ۹۳ متغیر برای هر ایستگاه ذخیره می‌شود. متغیرها شامل آماره‌های پایه (mean, std, skewness, median, count)، بهترین توزیع (best\_dist)، پارامترهای هر توزیع و معیارهای اطلاعاتی (loglik, aicc, bic) هستند. فشرده‌سازی Zstd با سطح ۳ و chunking به‌صورت (366, 500) پیاده‌سازی شده است.

\subsection{موازی‌سازی}
سامانه از multiprocessing برای پردازش موازی بلوک‌ها پشتیبانی می‌کند. تعداد workerها از طریق پارامتر max\_workers در config.yaml قابل تنظیم است. در benchmarkهای گزارش‌شده از ۶ worker استفاده شده است.

% ============================================================
% ۵. نتایج
% ============================================================
\section{نتایج}

\subsection{عملکرد محاسباتی}
نتایج آزمون‌های عملکردی در جدول \ref{tab:performance} ارائه شده است. سامانه با سرعت ۷٫۵ ایستگاه در ثانیه و اوج حافظهٔ ۳٫۵۶ گیگابایت، قادر به پردازش کامل داده‌ها در حدود ۱۲٫۵ ساعت است. تعداد بالقوه‌ی عملیات برازش برای چهار مدل داده‌ی خام، با فرض اجرای هر چهار مدل برای هر ایستگاه و هر روز، برابر با ۳۳۸٬۶۲۷ × ۳۶۶ × ۴ ≈ ۴۹۵٬۷۵۰٬۰۰۰ است. حدود ۸.۴\% از موارد به دلیل داده‌ی ناکافی یا خطای عددی از پردازش حذف شده‌اند.

\begin{table}[h]
\centering
\caption{نتایج آزمون عملکردی سامانه}
\label{tab:performance}
\begin{tabular}{l r}
\toprule
\textbf{معیار} & \textbf{مقدار} \\
\midrule
تعداد کل ایستگاه‌ها & ۳۳۸٬۶۲۷ \\
تعداد روزهای سال & ۳۶۶ \\
تعداد سال‌های دوره & ۳۱ \\
تعداد توزیع‌های برازش‌شده (داده‌ی خام) & ۴ \\
تعداد برازش‌های بالقوه (داده‌ی خام) & ۴۹۵٬۷۵۰٬۰۰۰ \\
تعداد برازش‌های معتبر & حدود ۴۵۴٬۵۰۰٬۰۰۰ (حدود ۹۱.۶\%) \\
سرعت پردازش & ۷٫۵ ایستگاه/ثانیه \\
زمان کل پردازش & ۱۲٫۵ ساعت \\
حافظۀ مصرفی (اوج) & ۳٫۵۶ گیگابایت \\
کاهش زمان I/O با کش & ۶۰\% \\
\bottomrule
\end{tabular}
\end{table}

\begin{table}[h]
\centering
\caption{مقایسهٔ زمان پردازش با و بدون کش برای بلوک‌های با اندازه‌های مختلف}
\label{tab:cache}
\begin{tabular}{l c c c}
\toprule
\textbf{اندازهٔ بلوک} & \textbf{بدون کش (ثانیه)} & \textbf{با کش (ثانیه)} & \textbf{کاهش (\%)} \\
\midrule
۱۰۰۰ & ۱۲۰ & ۴۸ & ۶۰ \\
۲۰۰۰ & ۲۳۰ & ۹۰ & ۶۱ \\
۵۰۰۰ & ۵۵۰ & ۲۱۵ & ۶۱ \\
\bottomrule
\end{tabular}
\end{table}

\subsection{اعتبارسنجی شناسایی مدل و بازتوانی پارامترها}
برای ارزیابی صحت روش انتخاب مدل، داده‌های مصنوعی با توزیع‌های معلوم و پارامترهای مشخص تولید شدند. پارامترهای توزیع‌های مولد به‌صورت زیر در نظر گرفته شدند: Normal با $\mu=0,\sigma=1$؛ Skew-Normal با $\xi=0,\omega=1,\alpha=2$؛ مخلوط گاوسی دومؤلفه‌ای با $w_1=0.4,\mu_1=-2,\sigma_1=1,\mu_2=2,\sigma_2=1$؛ و گامای انتقال‌یافته با $\alpha=4,\beta=1,\gamma=0$. داده‌ها با استفاده از seed ثابت ۴۲ و در ۱۰۰۰ تکرار مستقل برای هر سناریو تولید شدند. در هر تکرار، همان قاعده‌ی انتخاب مورد استفاده در داده‌های واقعی اعمال شد؛ یعنی کمینه‌ی AICc به‌عنوان معیار اصلی در نظر گرفته شد و در صورت وجود اختلاف کمتر از ۲ واحد بین مدل‌های رقیب، مدل با تعداد پارامتر کمتر انتخاب شد. جدول \ref{tab:validation} درصد انتخاب صحیح مدل را به‌همراه فاصله‌ی اطمینان ۹۵\% (با روش نرمال تقریبی) نشان می‌دهد. اعتبارسنجی در این بخش صرفاً مربوط به چهار مدل مورد استفاده در انتخاب توزیع برای داده‌های خام است؛ از آنجا که GEV در این مطالعه در یک گردش‌کار جداگانه و بر مبنای بیشینه‌های بلوک‌های سالانه برازش می‌شود، ارزیابی آن در جدول حاضر وارد نشده است.

\begin{table}[h]
\centering
\caption{عملکرد انتخاب مدل در داده‌های مصنوعی (۱۰۰۰ تکرار). مقادیر شامل درصد انتخاب صحیح و فاصله‌ی اطمینان ۹۵\% هستند.}
\label{tab:validation}
\begin{tabular}{l c c c}
\toprule
\textbf{توزیع مولد} & \textbf{$n=30$} & \textbf{$n=50$} & \textbf{$n=155$} \\
\midrule
نرمال & ۸۷.۰\% [۸۴.۹, ۸۹.۱] & ۹۳.۰\% [۹۱.۴, ۹۴.۶] & ۹۷.۰\% [۹۶.۰, ۹۸.۰] \\
نرمال چوله & ۷۹.۰\% [۷۶.۴, ۸۱.۶] & ۸۶.۰\% [۸۳.۸, ۸۸.۲] & ۹۲.۰\% [۹۰.۳, ۹۳.۷] \\
مخلوط گاوسی دومؤلفه‌ای & ۷۱.۰\% [۶۸.۲, ۷۳.۸] & ۸۰.۰\% [۷۷.۵, ۸۲.۵] & ۸۸.۰\% [۸۶.۰, ۹۰.۰] \\
گامای انتقال‌یافته & ۷۶.۰\% [۷۳.۳, ۷۸.۷] & ۸۴.۰\% [۸۱.۷, ۸۶.۳] & ۹۰.۰\% [۸۸.۱, ۹۱.۹] \\
\bottomrule
\end{tabular}
\end{table}

همچنین بازتوانی پارامترها برای مدل مخلوط گاوسی دومؤلفه‌ای با پارامترهای واقعی $w_1=0.4,\mu_1=-2,\sigma_1=1,\mu_2=2,\sigma_2=1$ در $n=155$ نشان داد که خطای نسبی میانگین (میانگین قدرمطلق خطای نسبی در ۱۰۰۰ تکرار) برای $\mu$ حدود ۲.۱\%، برای $\sigma$ حدود ۴.۳\% و برای $w$ حدود ۳.۷\% است.

\subsection{توزیع تابع‌های برازش‌شده برای داده‌های خام}
تحلیل توزیع تابع‌های برازش‌شده برای متغیر Tmax در جدول \ref{tab:selection} ارائه شده است. توزیع نرمال با ۴۲.۳\% بیشترین سهم را دارد، اما ۵۷.۷\% انتخاب‌ها به مدل‌های غیرنرمال اختصاص دارد. تحلیل ΔAICc نشان می‌دهد که در حدود ۶۸\% موارد، اختلاف AICc بین بهترین و دومین مدل بیشتر از ۲ است که بیانگر حمایت نسبتاً قوی از مدل منتخب است.

\begin{table}[h]
\centering
\caption{توزیع تابع‌های برازش‌شده برای متغیر Tmax (داده‌ی خام پنجره)}
\label{tab:selection}
\begin{tabular}{l r}
\toprule
\textbf{مدل} & \textbf{سهم انتخاب (\%)} \\
\midrule
نرمال & ۴۲.۳ \\
نرمال چوله & ۲۸.۷ \\
مخلوط گاوسی دومؤلفه‌ای & ۱۵.۲ \\
گامای انتقال‌یافته & ۵.۴ \\
بدون برازش معتبر & ۸.۴ \\
\bottomrule
\end{tabular}
\end{table}

\subsection{تحلیل دوکوهانه}
مدل مخلوط گاوسی دومؤلفه‌ای در ۱۵.۲\% از موارد انتخاب شده است. شاخص Ashman's D برای این موارد به‌طور میانگین برابر ۲.۳ بوده که نشان‌دهندهٔ جدایش واضح دو مؤلفه است. ضریب همپوشانی (OVL) میانگین ۰.۲۴ و شاخص جدایش وزنی (β) میانگین ۱.۸ محاسبه شد. تحلیل حساسیت با تغییر مقداردهی اولیهٔ EM نشان داد که در بیش از ۹۴\% موارد، مدل نهایی یکسان باقی می‌ماند که بیانگر پایداری الگوریتم است.

\subsection{نتایج تحلیل GEV برای مقادیر حدی}
تحلیل GEV بر روی بیشینه‌های سالانه‌ی پنجره برای متغیر Tmax نشان داد که میانگین برآورد نقطه‌ای پارامتر شکل ($\xi$) در سطح شبکه برابر با $0.12$ با انحراف معیار $0.04$ (پراکندگی بین ایستگاه‌ها) است. بررسی توزیع $\xi$ نشان داد که حدود ۶۲\% از ایستگاه‌ها دارای $\hat\xi>0$ هستند که به‌عنوان شواهد اکتشافی از گرایش به رفتار دم سنگین (Fréchet) در بخش قابل‌توجهی از شبکه در نظر گرفته می‌شود. با این حال، عدم‌قطعیت برآوردها قابل‌توجه است و این نتایج باید به‌عنوان شواهد اکتشافی تفسیر شوند، نه اثبات قطعی دم سنگین برای ایستگاه‌های منفرد. دوره‌های بازگشت محاسبه‌شده برای آستانه‌ی $T=10$ سال، $T=20$ سال و $T=50$ سال به‌ترتیب نشان‌دهنده‌ی افزایش قابل‌توجه دما در رویدادهای فرین نسبت به میانگین هستند. این نتایج به‌عنوان تحلیل مکمل ارائه می‌شوند و با نتایج انتخاب مدل برای داده‌های خام قابل‌مقایسه مستقیم نیستند.

\textbf{توجه:} با توجه به حجم نمونه‌ی محدود (۳۱ مقدار سالانه) و حساسیت برآوردها به پارامتر شکل، توصیه می‌شود دوره‌های بازگشت بزرگ‌تر از ۲۰ سال با احتیاط و به‌عنوان برون‌یابی تفسیر شوند. در نسخه‌ی حاضر، عدم‌قطعیت دوره‌های بازگشت به‌صورت کامل کمی نشده است و این موضوع به‌عنوان یکی از محدودیت‌های اصلی تحلیل GEV باقی می‌ماند.

% ============================================================
% ۶. بحث
% ============================================================
\section{بحث}

\subsection{ناهمگنی آماری داده‌های اقلیمی}
نتایج انتخاب مدل برای Tmax نشان می‌دهد که تنها ۴۲.۳\% از برازش‌ها به توزیع نرمال اختصاص دارد. این یعنی در بیش از نیمی از موارد، یک توزیع غیرنرمال (چوله، مخلوط گاوسی یا گامای انتقال‌یافته) بر اساس معیار AICc ترجیح داده شده است. این الگو تأکید می‌کند که فرض یک توزیع ثابت برای کل شبکه‌ی ایستگاهی، ناهمگنی آماری داده‌ها را نادیده می‌گیرد. مدل مخلوط گاوسی دومؤلفه‌ای با ۱۵.۲\% انتخاب، نشان‌دهندهٔ حمایت آماری از ساختار دو مؤلفه‌ای در بخش قابل‌توجهی از داده‌هاست، هرچند تفسیر فیزیکی آن نیازمند تحلیل مستقل است.

\subsection{تفسیر مدل مخلوط گاوسی دومؤلفه‌ای}
انتخاب مدل مخلوط گاوسی دومؤلفه‌ای باید با احتیاط تفسیر شود. وجود دو مؤلفه در مدل آماری الزاماً به معنای وجود دو رژیم فیزیکی مستقل نیست. این می‌تواند ناشی از ترکیب داده‌های فصلی، تغییرات ارتفاعی، ناهمگونی مکانی، یا حتی مصنوعات داده باشد. با این حال، شاخص‌های Ashman's D ($\approx 2.3$)، OVL ($\approx 0.24$) و شاخص جدایش وزنی ($\approx 1.8$) نشان می‌دهند که جدایش دو مؤلفه در بسیاری از موارد قابل‌توجه است و ساختار دومؤلفه‌ای صرفاً یک مصنوع آماری نیست. برای تفسیر اقلیمی، نیاز به تحلیل مستقل با متغیرهای فیزیکی مانند ارتفاع، عرض جغرافیایی و اقلیم منطقه‌ای وجود دارد که به‌عنوان کار آینده پیشنهاد می‌شود.

\subsection{معنای محاسباتی معماری}
ترکیب پردازش بلوکی، کش‌سازی هوشمند و سیستم checkpoint یک pipeline مقاوم و کارآمد ایجاد کرده است. کاهش ۶۰\% زمان مرتبط با ورودی/خروجی با کش، تأثیر مستقیمی بر زمان کل اجرا دارد. پردازش بلوکی با block\_size=۲۰۰۰ تعادل مناسبی بین مصرف حافظه (۳٫۵۶ گیگابایت) و زمان پردازش (۱۲٫۵ ساعت) ایجاد کرده است. سیستم checkpoint نیز امکان اجرای طولانی‌مدت را در محیط‌های عملیاتی فراهم می‌کند. اعتبارسنجی با داده مصنوعی نشان داد که سامانه با دقت بالایی قادر به تشخیص توزیع مولد و بازتوانی پارامترهاست.

\subsection{محدودیت‌ها}
چارچوب حاضر دارای محدودیت‌های زیر است:
\begin{enumerate}
\item استقلال کامل مشاهدات در پنجره برای برازش توزیع‌های مقطعی فرض شده است؛ در حالی که داده‌های پنجره‌ای دارای وابستگی زمانی هستند. این وابستگی در مدل‌سازی فعلی لحاظ نشده و به‌عنوان یک تقریب محاسباتی در نظر گرفته شده است.
\item انتخاب مدل مبتنی بر AICc مدل برتر را در میان مجموعه‌ی محدودی از چهار توزیع نامزد مشخص می‌کند، اما تضمینی برای اینکه توزیع مولد واقعی در این مجموعه قرار دارد، وجود ندارد.
\item الگوریتم EM ممکن است به مقداردهی اولیه حساس باشد و در برخی موارد به بهینه‌های محلی همگرا شود، هرچند استفاده از سه initialization مختلف این ریسک را کاهش داده است.
\item تحلیل GEV بر روی ۳۱ مقدار بیشینه‌ی سالانه‌ی پنجره انجام شده است که حجم نمونه‌ی نسبتاً کوچکی است و عدم‌قطعیت برآوردها را افزایش می‌دهد.
\item مدل مخلوط گاوسی دومؤلفه‌ای صرفاً شواهد آماری از وجود ساختار دو مؤلفه‌ای ارائه می‌کند و به‌تنهایی علت فیزیکی آن را اثبات نمی‌کند.
\item وابستگی فضایی داده‌ها در مدل‌سازی فعلی لحاظ نشده است.
\item از آنجا که تعداد بسیار زیادی برازش مستقل در سطح ایستگاه-روز انجام می‌شود، درصدهای انتخاب مدل در سطح شبکه ممکن است تحت تأثیر ناهمگنی حجم نمونه و ساختار وابستگی فضایی و زمانی قرار گیرند. در نسخه‌ی حاضر، اصلاح چندگانگی آزمون‌ها انجام نشده است، زیرا AICc به‌عنوان معیار انتخاب مدل و نه آزمون فرضیه استفاده شده است.
\item عدم‌قطعیت دوره‌های بازگشت GEV ناشی از حجم نمونه‌ی کوچک و برون‌یابی، در این مطالعه به‌صورت کامل کمی نشده است.
\end{enumerate}

% ============================================================
% ۷. بازتولیدپذیری و کنترل کیفیت
% ============================================================
\section{بازتولیدپذیری و کنترل کیفیت}

برای تسهیل بازتولیدپذیری نتایج، موارد زیر در مخزن GitHub ثبت شده‌اند: شناسه commit، نسخه‌ی پایتون (۳٫۱۱)، نسخه‌ی کتابخانه‌های اصلی (NumPy 1.26، SciPy 1.11، Xarray 2024.2، Zarr 2.17، Numba 0.59)، مشخصات سخت‌افزار (CPU Intel Xeon Gold 6254، ۱۲ هسته فیزیکی، ۲۴ رشته، ۶۴ گیگابایت RAM، NVMe SSD، OS Ubuntu 22.04 LTS)، تعداد workerهای موازی (۶)، اندازه‌ی بلوک (۲۰۰۰)، پیکربندی کش (فعال، حداکثر ۲۰ گیگابایت) و seed تصادفی (۴۲). نسخهٔ بایگانی‌شده‌ی پروژه متعاقباً در Zenodo با شناسه‌ی دائمی (DOI) منتشر خواهد شد و لینک آن در مخزن به‌روزرسانی می‌گردد.

سیستم کنترل کیفیت شامل پرچم‌های زیر است که برای هر برازش محاسبه می‌شوند:
\begin{itemize}
\item \textbf{PASS}: برازش معتبر
\item \textbf{LOW\_SAMPLE}: تعداد مشاهدات کمتر از آستانهٔ مدل
\item \textbf{NO\_CONVERGENCE}: همگرایی الگوریتم برقرار نشد
\item \textbf{OUTLIER\_CANDIDATE}: وجود داده‌های پرت (بیشتر از ۴$\sigma$)
\item \textbf{HIGH\_AICC}: AICc بالاتر از آستانه (۱۰۰۰)
\item \textbf{BAD\_SKEW}: چولگی خارج از محدوده قابل‌قبول
\item \textbf{NAN\_INPUT}: داده‌های ورودی شامل NaN
\end{itemize}
پرچم‌ها به‌صورت بیت‌ماسک ذخیره می‌شوند تا امکان ثبت همزمان چندین وضعیت فراهم شود.

% ============================================================
% ۸. نتیجه‌گیری و کارهای آینده
% ============================================================
\section{نتیجه‌گیری و کارهای آینده}

در این پژوهش، موتور پردازش اقلیم‌شناسی به‌عنوان یک چارچوب کارآمد و مقیاس‌پذیر برای برازش توزیع‌های احتمالاتی بر داده‌های سری‌زمانی اقلیمی معرفی شد. یافته‌های آماری نشان می‌دهد که توزیع نرمال با ۴۲.۳\% بیشترین سهم را دارد، اما ۵۷.۷\% انتخاب‌ها به مدل‌های غیرنرمال اختصاص دارد. مدل مخلوط گاوسی دومؤلفه‌ای با ۱۵.۲\% انتخاب، نشان‌دهندهٔ حمایت آماری از ساختار دو مؤلفه‌ای در بخش قابل‌توجهی از داده‌هاست، هرچند تفسیر فیزیکی آن نیازمند تحلیل مستقل است. تحلیل جداگانه‌ی GEV نیز گرایش به رفتار دم سنگین را در بخش قابل‌توجهی از ایستگاه‌ها نشان داد، هرچند عدم‌قطعیت برآوردها قابل‌توجه است. نتایج آزمون‌های عملکردی نشان داد که سامانه با سرعت ۷٫۵ ایستگاه در ثانیه و مصرف حافظهٔ کمتر از ۴ گیگابایت، قادر به پردازش کامل داده‌های یک دورهٔ ۳۱ ساله در کمتر از ۱۳ ساعت است. سیستم کش‌سازی هوشمند با کاهش ۶۰\% زمان مرتبط با ورودی/خروجی، نقش کلیدی در این کارایی دارد. اعتبارسنجی با داده‌های مصنوعی نشان داد که سامانه در بازتوانی توزیع‌های مولد و پارامترهای آن‌ها عملکرد مطلوبی دارد.

پیشنهاد می‌شود در پژوهش‌های آینده: (۱) توزیع‌های بیشتری مانند Johnson SU، Wakeby و توزیع‌های شرطی به سامانه افزوده شوند، (۲) قابلیت‌های پردازش موازی با استفاده از Dask و Ray توسعه یابد، (۳) پردازش داده‌های شبکه‌ای (gridded) و خروجی در فرمت NetCDF با استاندارد CF پیاده‌سازی شود، (۴) تحلیل فضایی انتخاب مدل‌ها و ارتباط آن با عوامل جغرافیایی و اقلیمی انجام شود، (۵) اعتبارسنجی پیش‌بینانه با cross-validation و آزمون‌های goodness-of-fit افزوده شود، و (۶) عدم‌قطعیت پارامترها با bootstrap و عدم‌قطعیت مدل با model averaging کمی شود.

% ============================================================
% تشکر و قدردانی
% ============================================================
\section*{تشکر و قدردانی}
این پژوهش با حمایت مرکز ملی خشکسازی و سازمان هواشناسی کشور انجام شده است. نویسنده از همکاران خود در اداره پایش اقلیمی به‌خاطر همکاری و ارائهٔ بازخوردهای سازنده تشکر می‌کند. کد منبع این پروژه به‌صورت متن‌باز در مخزن GitHub به آدرس \href{https://github.com/AminFazlKazemi/ClimateProcessingEngine}{github.com/AminFazlKazemi/ClimateProcessingEngine} در دسترس است. نسخهٔ بایگانی‌شده‌ی پروژه متعاقباً در Zenodo با شناسه‌ی دائمی (DOI) منتشر خواهد شد و لینک آن در مخزن به‌روزرسانی می‌گردد.

% ============================================================
% منابع
% ============================================================
\begin{thebibliography}{99}

\bibitem{wilks2011}
Wilks, D. S. (2011). \textit{Statistical Methods in the Atmospheric Sciences} (3rd ed.). Academic Press.

\bibitem{katz2002}
Katz, R. W., Parlange, M. B., \& Naveau, P. (2002). Statistics of extremes in hydrology. \textit{Advances in Water Resources}, 25(8-12), 1287–1304.

\bibitem{coles2001}
Coles, S. (2001). \textit{An Introduction to Statistical Modeling of Extreme Values}. Springer.

\bibitem{dempster1977}
Dempster, A. P., Laird, N. M., \& Rubin, D. B. (1977). Maximum likelihood from incomplete data via the EM algorithm. \textit{Journal of the Royal Statistical Society: Series B}, 39(1), 1–38.

\bibitem{hosking1990}
Hosking, J. R. M. (1990). L-moments: Analysis and estimation of distributions using linear combinations of order statistics. \textit{Journal of the Royal Statistical Society: Series B}, 52(1), 105–124.

\bibitem{mclachlan2000}
McLachlan, G. J., \& Peel, D. (2000). \textit{Finite Mixture Models}. Wiley.

\bibitem{azzalini1985}
Azzalini, A. (1985). A class of distributions which includes the normal ones. \textit{Scandinavian Journal of Statistics}, 12(2), 171–178.

\bibitem{burnham2002}
Burnham, K. P., \& Anderson, D. R. (2002). \textit{Model Selection and Multimodel Inference: A Practical Information-Theoretic Approach} (2nd ed.). Springer.

\bibitem{github}
Fazl Kazemi, A. (2024). ClimateProcessingEngine: A scalable framework for probability distribution fitting in climate time series. GitHub repository. \url{https://github.com/AminFazlKazemi/ClimateProcessingEngine}

\bibitem{xarray}
Hoyer, S., \& Hamman, J. (2017). Xarray: N-D labeled arrays and datasets in Python. \textit{Journal of Open Research Software}, 5(1), 10.

\bibitem{zarr}
Miles, A., et al. (2020). Zarr: A portable chunked, compressed, N-dimensional array format for Python. \textit{Journal of Open Source Software}, 5(50), 2199.

\bibitem{numba}
Lam, S. K., Pitrou, A., \& Seibert, S. (2015). Numba: A LLVM-based Python JIT compiler. \textit{Proceedings of the Second Workshop on the LLVM Compiler Infrastructure in HPC}, 1–6.

\end{thebibliography}

% ============================================================
% ضمیمه
% ============================================================
\onecolumn
\section*{ضمیمه: الگوریتم‌ها و کدهای اصلی}

\subsection*{الگوریتم ۱: پردازش بلوکی و انتخاب مدل برای داده‌های خام}
\begin{algorithm}[h]
\caption{پردازش بلوکی و انتخاب مدل (داده‌ی خام)}
\begin{algorithmic}[1]
\State \textbf{ورودی:} داده‌های اقلیمی (N\_stations، N\_years، N\_days، N\_vars)
\State \textbf{خروجی:} خروجی Zarr با متغیرهای ۹۳گانه
\State
\State بارگذاری پیکربندی و مقداردهی اولیه Zarr
\State بررسی checkpoint و تعیین نقطه‌ی شروع
\For{هر بلوک از ایستگاه‌ها (block\_size)}
	\State بارگذاری داده‌های بلوک (از کش یا Zarr)
	\For{هر ایستگاه در بلوک}
		\For{هر روز از سال}
			\State استخراج پنجره‌ی ۵ روزه
			\State حذف مقادیر نامعتبر (NaN)
			\State محاسبه‌ی تعداد مشاهدات معتبر ($n$)
			\For{هر مدل توزیعی $m \in \mathcal{M}_{raw}$ (Normal, Skew-Normal, GMM, Shifted Gamma)}
				\If{$n \ge n_{\min}^{(m)}$}
					\State برازش مدل و محاسبهٔ log-likelihood
					\State محاسبهٔ AICc و BIC
					\If{برازش ناموفق}
						\State ثبت پرچم کیفیت مناسب
					\EndIf
				\Else
					\State ثبت \texttt{LOW\_SAMPLE}
				\EndIf
			\EndFor
			\If{حداقل یک برازش موفق وجود دارد}
				\State $\Delta AICc_i = AICc_i - \min(AICc)$
				\State انتخاب مدل با کمترین AICc؛ در صورت $\Delta AICc < 2$، انتخاب مدل با کمترین $k$
				\State ذخیره پارامترها، معیارهای اطلاعاتی و پرچم کیفیت
			\Else
				\State ذخیرهٔ \texttt{NO\_VALID\_FIT}
			\EndIf
		\EndFor
	\EndFor
	\State نوشتن نتایج بلوک در Zarr
	\State به‌روزرسانی checkpoint
\EndFor
\State نهایی‌سازی متادیتای Zarr
\State بازگرداندن مسیر خروجی
\end{algorithmic}
\end{algorithm}

\subsection*{الگوریتم ۲: تحلیل جداگانه‌ی GEV برای مقادیر حدی}
\begin{algorithm}[h]
\caption{تحلیل GEV بر روی بیشینه‌های سالانه‌ی پنجره حول روز تقویمی}
\begin{algorithmic}[1]
\State \textbf{ورودی:} داده‌های اقلیمی، ایستگاه، روز هدف ($d$)
\State \textbf{خروجی:} پارامترهای GEV و دوره‌های بازگشت
\State
\For{هر سال $y \in \{1369,\dots,1399\}$}
	\State استخراج پنجره‌ی ۵ روزه: $X_{y,d-2}, \dots, X_{y,d+2}$
	\State محاسبه‌ی بیشینه: $M_y(d) = \max(X_{y,d-2}, \dots, X_{y,d+2})$
\EndFor
\State تشکیل دنباله‌ی $\{M_{1369}(d), \dots, M_{1399}(d)\}$
\If{تعداد مقادیر معتبر $\ge 10$}
	\State برازش GEV با \texttt{genextreme.fit}
	\State محاسبه‌ی دوره‌های بازگشت $z_T$ برای $T \in \{2,5,10,20,50\}$
\Else
	\State ثبت \texttt{LOW\_SAMPLE\_GEV}
\EndIf
\State بازگرداندن پارامترها و دوره‌های بازگشت
\end{algorithmic}
\end{algorithm}

\end{document}